\documentclass[a6paper, 11pt, parskip=half]{scrartcl}

\usepackage[dvipsnames]{xcolor}
\usepackage{tikz}

\usepackage{ragged2e}
% Minimize unwanted hyphenation
\tolerance=1
\emergencystretch=\maxdimen
\hyphenpenalty=1
\hbadness=10000

\usepackage{fontspec}
\setmainfont{Cochin}
%\setmainfont{Qwigley}

\setkomafont{section}{\setmainfont{Able New}\LARGE}
\setkomafont{subsection}{\setmainfont{Able New}\Large}
\setkomafont{subsubsection}{\setmainfont{Able New}\large}

% Adjust spacing before and after section headings
\RedeclareSectionCommand[
  runin=false,
  beforeskip=1.0\baselineskip,
  afterskip=-0.0\baselineskip
]{section}

% Adjust spacing before and after subsection headings
\RedeclareSectionCommand[
  runin=false,
  beforeskip=1.0\baselineskip,
  afterskip=-0.0\baselineskip
]{subsection}

% Adjust spacing before and after subsubsection headings
\RedeclareSectionCommand[
  runin=false,
  beforeskip=1.0\baselineskip,
  afterskip=-0.0\baselineskip
]{subsubsection}


\usepackage{enumitem}
\setlist[description]{labelindent=0pt, labelsep=\widthof{ }, leftmargin=\widthof{\textbf{License: }}, font=\setmainfont{Cochin}\bfseries}

\usepackage[hang,flushmargin]{footmisc}
\newcommand\blfootnote[1]{%
  \begingroup
  \renewcommand\thefootnote{}\footnote{#1}%
  \addtocounter{footnote}{-1}%
  \endgroup
}
\renewcommand{\thefootnote}{\fnsymbol{footnote}}
\renewcommand{\footnoterule}{%
  \kern -3pt
  \hrule width \textwidth height 0.5pt
  \kern 2pt
}

\usepackage[hidelinks]{hyperref}
\usepackage[type={CC}, version={4.0}, modifier={by-sa}]{doclicense} % Add text and icons for creative commons license
%\usepackage{array}

\raggedright
\pagestyle{empty}
\begin{document}

\begin{titlepage}
\enlargethispage{4.75\baselineskip}
\setmainfont[Scale=3.25]{Magic School One}
\Huge
\begin{center}
\vspace*{-1.0\baselineskip}
Staff

Meeting

\vfill
\includegraphics[scale=0.033]{Images/wizard_with_staff.png}
\vfill
\huge
%\setmainfont{Cochin}
\setmainfont{Bilbo Swash Caps}
A Socratic Worldbuilding Game\\
\setmainfont{Bilbo Swash Caps}
by Michael Purcell
\end{center}
\end{titlepage}
\thispagestyle{empty}
\enlargethispage{3.5\baselineskip}
\setmainfont{Bilbo Swash Caps}
\large
\noindent Esteemed Colleagues,

As you may be aware, several companies have recently developed magical devices commonly known simply as ``staffs''.
These devices allow their users to cast spells in new and powerful ways. 
A few pioneering students have already started using staffs.
Doubtless, many more will soon follow their example.

So, we should to try to understand what the impact of these staffs will be for us here at the school.
To that end, you are cordially invited to attend a mandatory meeting tomorrow to discuss the matter.
I have identified the following topics as being of particular interest:
\begin{enumerate}[nosep]
	\item Educational Outcomes
	\item Safety \& Logistics
	\item Academic Integrity
	\item Equal Opportunity
\end{enumerate}   

Please familiarize yourself with some aspects of how staffs work and what they can do ahead of time.
By sharing our expertise, we can all learn a bit more about this exciting and powerful new technology.

\hspace{4.5cm}\huge- M
\setmainfont{Cochin}
\normalsize

\newpage
\enlargethispage{3.5\baselineskip}

\section*{Overview}
This is a worldbuilding game. It is intended for groups of three to six players and can be played in about one hour.

You will assume the roles of teachers at a school of witchcraft and wizardry. Your students have recently started using staffs, a new kind of device that changes the ways in which they can interact with magic. You have gathered to discuss the situation and decide how you want to respond.

During the game, you will address a series of questions about how staffs work and how they will affect life at the school. You will invent the answers to these questions as you play. By doing so, you will describe the world in which your characters live, how magic works there, and how it will be affected by this new technology.

\newpage
\enlargethispage{3.5\baselineskip}

\section*{Characters}
You will portray a teacher who works at the school. To create your character,
\begin{enumerate}[nosep]
	\item Describe the subject that you teach.
	\item Describe your mannerisms and physical appearance.
	\item Describe two types of questions about staffs that you can answer.
	\item Describe one objective that you want to accomplish at the meeting.
\end{enumerate}
Introduce yourself to the other characters before the meeting begins.% You do not have to reveal your objective at this time.

\section*{Gameplay}
The game takes place over a series of five rounds. 
During each of the first four rounds, you will discuss one topic that the headmaster has identified as being of particular interest.
During the last round, you will discuss how to respond to the issues raised in previous rounds. 

\newpage
\enlargethispage{3.5\baselineskip}

\subsection*{Topics}
The headmaster has identified four topics of particular interest for discussion:
\begin{enumerate}[nosep]
	\item Educational Outcomes
	\item Safety \& Logistics
	\item Academic Integrity
	\item Equal Opportunity
\end{enumerate}
In each of the first four rounds, you will discuss one of these topics.

\subsection*{Facilitators}
One player should act as the facilitator in each round. Their job is to ensure that everyone has a fair chance to contribute and that the discussion stays focused on the topic at hand. In addition to their administrative duties, the facilitator may (and should) contribute to the discussion via their character as usual. 

\newpage
\enlargethispage{3.5\baselineskip}

\subsection*{Questions}
During the game, you will ask and answer questions both about staffs and about the game's setting.
You will invent the answers to these questions as they arise to describe the world that your characters inhabit.

\subsubsection*{Factual Questions}

Factual questions are questions about the nature of some part of the game's setting.
These questions are usually stated using one of the standard interrogatives: who, what, where, when, why, and how.

Whoever is best suited to answer each factual question should do so.
Your answers should be consistent with any details about the setting that have been previously established.
Beyond that, however, you are free to invent any details you like as a part of your answers.

\newpage
\enlargethispage{3.5\baselineskip}

\subsubsection*{Socratic Questions}
Socratic questions are questions that encourage critical thinking.
These questions frequently arise as follow-up questions after a player establishes a new detail about the game's setting.

Socratic questions are often intended to do one or more of the following:
\begin{itemize}[nosep]
	\item Clarify concepts
	\item Challenge assumptions
	\item Probe evidence
	\item Discover alternative viewpoints
	\item Explore implications
\end{itemize}

After someone answers a factual question, you should use Socratic questions to help them flesh out their answer and explain how any new details they introduced interact with other details that had been previously established.

\newpage
\enlargethispage{3.5\baselineskip}

\subsection*{Action Items}
During the last round, your job is to decide what you are going to do next.
You should describe what you think needs to be done, what you can do yourselves, and what you need help with.

As in the previous rounds, you should use Socratic questions to help each other understand what you are proposing and why you think that your proposal describes a reasonable course of action.

\vfill
\begin{center}
\includegraphics[scale=0.09]{Images/druid_Staff.png}
\end{center}
\vfill
%\hrulefill
\begin{description}[nosep]
\item[\textbf{Design}:] Michael Purcell \\
\item[\textbf{Contact}:] \href{mailto:sm.ttrpg@gmail.com}{sm.ttrpg@gmail.com}\\
\item[\textbf{License}:] \doclicenseText%This work is licensed under a ``CC BY 4.0'' license.
\end{description}
%\begin{tabular}{@{}m{\columnwidth-\widthof{\Huge{\doclicenseIcon}}-0.5cm}@{\hspace{0.05cm}}m{\widthof{\Huge{\doclicenseIcon}}}@{}}
%{\textbf{License}: This work is licensed\newline under a ``CC BY 4.0'' license.} & \Huge{\doclicenseIcon}\\
%\end{tabular}

\end{document}